%% ===========================================================================
%%  frontmatter/daftar-singkatan.tex — Daftar Singkatan dan Akronim
%%  Spasi 1,5, urut abjad. Format entri: \singkatan{SINGKATAN}{Keterangan}
%% ===========================================================================
\babfrontmatter{Daftar Singkatan dan Akronim}

\begin{daftarsingkatan}
\singkatan{AI}{\asing{Artificial Intelligence} (kecerdasan buatan)}
\singkatan{API}{\asing{Application Programming Interface}}
\singkatan{BMW}{\asing{Broadband Multimedia Wireless}}
\singkatan{GitOps}{Praktik pengelolaan infrastruktur melalui repositori Git}
\singkatan{GPU}{\asing{Graphics Processing Unit}}
\singkatan{IEEE}{\asing{Institute of Electrical and Electronics Engineers}}
\singkatan{JSON}{\asing{JavaScript Object Notation}}
\singkatan{LLM}{\asing{Large Language Model} (model bahasa besar)}
\singkatan{MCP}{\asing{Model Context Protocol}}
\singkatan{NIDN}{Nomor Induk Dosen Nasional}
\singkatan{NIM}{Nomor Induk Mahasiswa}
\singkatan{NTUST}{National Taiwan University of Science and Technology}
\singkatan{O-Cloud}{\asing{Open Cloud} pada arsitektur \asing{O-RAN}}
\singkatan{PRD}{\asing{Product Requirements Document}}
\singkatan{QR}{\asing{Quick Response}}
\singkatan{RAG}{\asing{Retrieval-Augmented Generation}}
\singkatan{SAD}{\asing{System Architecture Design}}
\singkatan{SOTN}{\asing{School of Technopreneur Nusantara}}
\singkatan{SQLite}{\asing{Structured Query Language Lite}}
\singkatan{STMIK}{Sekolah Tinggi Manajemen Informatika dan Komputer}
\singkatan{TEEP}{\asing{Taiwan Education Experience Program}}
\singkatan{UI}{\asing{User Interface} (antarmuka pengguna)}
\end{daftarsingkatan}

